\section{实验}

本节利用黑潮区域六年的高分辨率再分析数据（第\ref{sec:scope}节另加湾流区与阿古拉斯回流区），依次回答六个问题：（1）3天前的1°粗图错在哪里，其误差是否具备可由大数据标定的物理结构？（2）单纯增加计算量能否消除选边错误，环境歧义有多普遍？（3）五类策略中哪种能真正降低航时损失？（4）损失集中在哪些任务上，最坏后悔为何优于集合翻转比例，何时值得去观测？（5）与随机合理场相比，反例搜索在算力上省下了什么？（6）当海流图变好、航行器变慢、更换误差模型、换至其他强流海区或改带真实卫星海流图时，方法的适用范围如何？

\subsection{实验设置与对比基线}
所有误差模态、标定半径（$r=1.5$）、数值误差容限（$\tau_{\text{num}}=3.2$小时）及门控阈值（$\theta=8$小时）均在2019--2021年开发年份（710个任务）上标定，原样用于2022--2024年测试年份（623个任务）评价。我们将14种对比方法分为五组（见图\ref{fig:main}）：（1）同一张粗图上加算（地图+粗算、地图+标准、照地图走（精算））；（2）不取新信息而改变决策准则（按期望时间、保守（最坏后悔最小））；（3）就地取新信息（原地停留测、沿途测30千米、沿途测60千米、绕行定点测）；（4）上浮取新图（总是上浮取当天1°图、总是上浮取当天0.5°图，每次计1小时）；（5）反例检验分流（门控后取保守、门控后上浮取图、门控后按期望时间）。显著性通过与照地图走（精算）配对差值的按天整组重抽样（2000次）95\%置信区间检验。

\subsection{旧粗图错在哪里：黑潮误差的空间结构与尺度}
3天前的1°粗图频繁误导决策，根本原因在于丢失了黑潮主轴及伴生涡旋的高幅值、长相关尺度切变，而这一结构可由多年历史数据稳定标定。开发年份的误差统计（图\ref{fig:data}）表明，黑潮强流带自台湾东北蜿蜒贯穿东海陆架坡折，3天前1°粗图的误差并非均匀白噪声，而是沿黑潮主轴形成高误差带，全区典型误差中位数为0.17米/秒，误差空间相关长度达76千米；相比之下，当天0.5°图的典型误差中位数降至0.07米/秒，相关长度缩短至18千米。76千米的误差相关长度意味着，在跨度百余千米的绕障任务中，一个未被粗图解析的中尺度误差模态足以笼罩整个左侧或右侧通道，使整条备选航线同时遭受逆流或顺流偏置，这构成了用低维流函数模态构造合理集合的物理依据。

\begin{figure}[htbp]
\centering
\includegraphics[width=\linewidth]{fig_route_data}
\caption{台湾以东至东海研究海区的流场特征与粗图误差标定。(a) 研究海区（23°--31°N、121°--129°E）在2019--2021年开发年份的10米平均海流场，黑潮主轴清晰可见，橙色圆点标示测试任务禁航区中心位置；(b) 3天前1°粗图的典型误差幅度分布（米/秒），最大误差带紧密贴合黑潮强流轴；(c) 海流误差相关系数随两点空间距离的衰减曲线，当天0.5°图的误差相关长度约为18千米，而3天前1°粗图达76千米。旧粗图的误差具有鲜明的空间结构与中尺度相关长度，可由多年海洋大数据可靠标定，构成了构造“合理海流”的物理依据。}
\label{fig:data}
\end{figure}

\subsection{计算消除不了环境歧义：算力与合理反例的实证分析}
在同一张3天前1°粗图上大幅增加计算量只能消除少部分数值歧义，而在包住真实海洋的合理集合内，几乎每个任务都存在翻转决策的合理反例。在测试年份中，将求解算力从粗算增加64倍至精算，改变了64个任务的选择，其中38个改对、26个改错，选错一侧的任务比例因此仅从21.7\%降至19.7\%，平均损失的变化在统计上也不显著；相反，若不加计算而换成当天0.5°图，选错率立即降至10.4\%（图\ref{fig:ambiguity}(a)）。多算无效的原因在于环境本身的多解性：在开发年份覆盖93\%真实时间差、测试年份覆盖94\%的标定半径 $r=1.5$ 处，99.7\%的测试任务都存在至少一个使另一侧更快的合理反例（图\ref{fig:ambiguity}(b)）。因此，在旧粗图下环境歧义是常态，仅问“是否存在反例”没有区分度，真正决定行动的是反例造成的最坏后悔有多大。

\begin{figure}[htbp]
\centering
\includegraphics[width=\linewidth]{fig_route_ambiguity}
\caption{机载计算量与环境信息对航路决策歧义的不同作用。(a) 测试年份选错绕行一侧的任务比例：在同一张3天前1°粗图上将计算量从粗算（21.7\%）增加至标准（20.5\%）与64倍精算（19.7\%），错误率仅下降约2个百分点，而换成当天0.5°图则直接减半至10.4\%（知道真相为0\%）；(b) 开发年份合理集合包住真实两侧时间差的任务比例与测试年份存在合理反例的任务比例随半径 $r$ 的变化，在标定半径 $r=1.5$ 处两者分别达93\%与99.7\%。机载计算只能消除少量的数值歧义；在足以包住真实海洋的合理集合里，几乎每个决定都存在合理反例，因此必须进一步度量反例导致的最坏损失大小。}
\label{fig:ambiguity}
\end{figure}

\subsection{主对比：五类响应策略的平均航时损失}
在五类响应策略的对比中，唯有基于合理集合控制尾部风险的保守准则与本文的门控后取保守策略显著降低了平均航时损失，而单纯加算、按期望时间选边与就地测流均未能显著改善结果（图\ref{fig:main}）。在测试年份上，照地图走（精算）的平均损失为0.98小时；改用32个合理场的按期望时间准则或门控后按期望时间，与精算地图的配对差置信区间均包含0，原因在于32个随机合理场数量太少、均值的抽样误差过大（见第\ref{sec:budget}节）。相反，门控后取保守将测试年份平均损失降至0.70小时，与照精算地图走的配对差为$-0.28$小时，95\%置信区间为$[-0.51, -0.07]$（不含0），相对损失降低28\%。门控后取保守之所以奏效，是因为在保守规则实际改选另一侧的41个任务中，有23个改对且平均每次省下10.1小时，18个改错但平均每次仅多花3.3小时——被避开的一侧往往具有很大的下行风险，改对的收益远大于改错的代价。同时，一律取保守的平均损失为0.72小时（与精算地图之差的置信区间亦不含0），门控后取保守与之相近：这表明不观测时的核心收益来自反例搜索给出的最坏后悔准则本身，而门控的作用在于锁定高风险任务，在调用有代价的上浮观测时将平均损失从总是上浮取0.5°图的1.22小时降至门控后上浮取图的0.78小时。

\begin{figure}[htbp]
\centering
\includegraphics[width=\linewidth]{fig_route_main}
\caption{测试年份623个强流任务上五类响应策略的平均航时损失（竖直虚线为照精算地图走的0.98小时，各策略的数值标在点旁）。五组从上至下依次为：同一张粗图上加算、不取新信息而改变决策准则、就地取新信息、总是上浮取新图（每次计1小时）、本文的反例检验门控分流。只有基于合理反例最坏后悔的保守与门控策略显著降低了损失，单纯加算与就地测流均难以奏效。}
\label{fig:main}
\end{figure}

\subsection{损失集中在哪里、何时值得观测}
航路决策的航时损失高度集中在少数尾部任务上，反例最坏后悔之所以有效，是因为它比集合翻转比例和地图时间差更能把大损失任务排在前面。若只论能否区分“是否选错且多花1小时以上”而不论多花多少，三种排序指标的区分能力其实相近；但考察排在前20\%的任务所捕获的地图总损失比例时，按反例最坏后悔达67\%，明显高于32个合理场翻转比例的49\%与地图两侧时间差的34\%（图\ref{fig:risk}(a)）。在门控阈值8小时下，被标记的47\%任务囊括了82\%的地图总损失，因为最坏后悔直接度量了不利流场造成的航时差，而非将小损失翻转与十几个小时的大损失翻转同等看待。

各类观测手段的成败则取决于观测视野与时间代价的权衡。就地原位测流均未能带来显著净收益（与精算地图之差的置信区间或显著为正、或包含0）：原地停留测耗去2小时却只测到起点一处；绕行定点测平均要先花15.4小时才能赶到前方测点，得不偿失；沿途测30千米后换边的任务仅占3.4\%（其中81\%换对），说明出发后前30千米测到的海流很少足以改变对两侧的判断，而沿途测60千米后换边的任务虽增至18.8\%，但由于已深入半程、折返代价过大，换边只有15\%是划算的。对于能更新全局流场的上浮取当天0.5°图（图\ref{fig:risk}(b)），当单次上浮代价为0.5小时时，门控后上浮取图的平均损失为0.55小时，优于门控后取保守的0.70小时；而当上浮代价升至1.0小时时，门控后上浮取图升至0.78小时，交点出现在0.5至1小时之间（近1小时），此后无需观测的门控后取保守成为更优选择。

\begin{figure}[htbp]
\centering
\includegraphics[width=\linewidth]{fig_route_risk}
\caption{黑潮测试年份的风险排序与上浮观测代价。(a) 按不同风险指标从高到低排序时，排在前面的任务累计占地图总损失的比例：排名前20\%的任务中，反例最坏后悔占67\%，明显高于32个合理场的翻转比例（49\%）与地图两侧时间差（34\%）；(b) 平均损失随单次上浮取当天0.5°图的时间代价（0--4小时）的变化：总是上浮随代价线性变差，门控后上浮在代价低于约1小时时最优，代价更高时门控后取保守最优。}
\label{fig:risk}
\end{figure}

\subsection{反例搜索省下了什么：与随机合理场的比较}
\label{sec:budget}
在决定选哪一侧时，反例搜索的优势在于以约1/29的求解次数达到与上千个随机合理场平均同等的降损效果。理论推导与逐任务核对表明，选最坏后悔更小的一侧等价于按合理集合给出的两侧时间差范围的中点选边；若两侧时间差随海流线性变化，中点即等于地图值，只有当一侧的最不利情形比另一侧严重得多时保守选择才会改选。测试年份实验（图\ref{fig:budget}(a)）显示，按期望时间选边随随机合理场数量增加而改善：用32个合理场时，其平均损失显著高于反例搜索后取保守（配对差的置信区间不含0），因为均值的抽样误差（中位数2.4小时）大于它要估计的修正量（中位数1.7小时），改选多由抽样噪声驱动；直到数百个合理场以上两者差异才不再显著，1024个合理场（2048次求解）的平均损失为0.71小时，与仅需约71次求解的反例搜索后取保守（0.72小时）相当，开发年份两者之差亦不显著。

在识别高风险任务时，排序有效的关键在于度量损失大小，而非翻转频率。测试年份中（图\ref{fig:budget}(b)），按随机合理场算出的期望后悔排在前20\%的任务，在32个场时已占地图总损失的65\%，接近反例搜索最坏后悔的67\%；翻转比例则即使用1024个随机合理场也只有44\%。因此，只要度量的是坚持地图一侧会多花多少航时，数十个随机合理场的期望后悔已能较好地排出危险任务；反例搜索为算力受限的航行器省下的，是做出正确选边所需的上千次航线求解。

\begin{figure}[htbp]
\centering
\includegraphics[width=\linewidth]{fig_route_budget}
\caption{随机合理场数量与反例搜索在选边决策和风险排序上的对比（测试年份）。(a) 平均航时损失随每个任务求解次数的变化：曲线为在 $n$ 个随机合理场上按期望时间选边（每次需 $2n$ 次求解），单点为反例搜索后取保守（约71次求解），水平虚线为照精算地图走（0.98小时）；(b) 按不同风险指标排名前20\%的任务所占地图总损失比例随求解次数的变化：两条曲线分别为同样随机合理场上的期望后悔与翻转比例，单点为反例搜索的最坏后悔（67\%），水平虚线为地图两侧时间差（34\%）。反例搜索以约1/29的求解次数达到了上千个随机合理场平均的选边效果，而风险排序的优劣取决于是否度量损失大小，而非样本数量。}
\label{fig:budget}
\end{figure}

\subsection{适用范围与消融分析}
\label{sec:scope}
表\ref{tab:scope}进一步检验了不同先验地图、航速、误差模型以及另外两个强流海区下的适用边界。第一，当航行器直接携带当天0.5°图时，存在合理反例的任务降至64.5\%，精算地图平均损失仅0.17小时，门控后取保守自动退化为绝大多数时候照地图走（损失同为0.17小时，配对差置信区间含0），不产生负向干扰。第二，当航行器对水速度降至0.75米/秒时，选错一侧的代价被显著放大，精算地图平均损失升至2.43小时；门控后取保守将损失降至1.80小时（但配对差置信区间包含0，统计上不显著），而由于选错代价变大、上浮相对更划算，门控后上浮取图将损失进一步降至1.32小时。第三，在消融实验中，若改用历史相似日的小尺度差构造合理集合，门控后取保守的平均损失为0.95小时，相比精算地图的0.99小时几乎失去纠错能力（配对差置信区间含0）。两组的反例搜索方法完全相同，差别只在“合理”如何定义：由多年数据标定误差大小与相关长度的合理集合，能够指出真正危险的一侧，而拼装历史相似日的小尺度差则不能。

第四，在北大西洋湾流海区（33°--41°N、75°--67°W）的重复实验中，“计算消除不了环境歧义”“最坏后悔能找出损失所在”与“门控后去观测”三条结论均得到复现。该实验采用同一再分析产品、深度、年份划分与流程，半径与阈值由湾流开发年份重新标定。湾流的旧粗图误差更大，照精算地图走的平均损失达3.56小时；加算64倍改变了22个任务的选择，其中8个改对、14个改错；门控后上浮取当天0.5°图将平均损失降至1.31小时（与精算地图之差的置信区间不含0），也优于总是上浮取图。

然而，不观测直接改选保守一侧在湾流未能复现：门控后取保守的平均损失为4.00小时，与精算地图之差的置信区间包含0。原因在于湾流测试任务中有8.7\%在真实海流下有一侧因顶流过强而不可通行（黑潮为1.0\%）；这类任务的合理集合内既有堵死左侧、也有堵死右侧的流场，两侧最坏后悔同时接近上限，选较小的一侧近于随机，共有5次改选误入真实不可通行的一侧。看过测试结果后所做的事后检验（非预先设定）表明，若仅在两侧最坏后悔之差超过数值误差容限时才改选，湾流的损失回到3.45小时，与照地图走无显著差别，而黑潮仍保留大部分收益。这提示：合理反例检验能可靠地判断粗图何时不可信，此时去观测是稳健的回应；不观测的保守改选可能只在两侧最坏情形差别分明时有效，下面在第三海区检验这一点。

第五，在阿古拉斯回流区（35°--43°S、14°--22°E）的第三海区独立检验中，加算无效、风险排序与门控后观测三条结论再次成立，而改进规则无害但无显著收益。该检验的计划在下载数据之前写定且未作改动，专门检验湾流实验后提出的改进规则（仅当两侧最坏后悔之差超过数值误差容限时才改选，不引入新参数）；半径与阈值由该区开发年份标定，测试年份只评价一次。预先写定的五项判定结果为：改进规则与精算地图之差为$-0.16$小时，置信区间含0；原门控后取保守同样不显著（该区16.2\%的任务有一侧不可通行）；门控后上浮取当天0.5°图将平均损失从6.87小时降至2.20小时，置信区间不含0；加算64倍后选错率不变（16个选择改对、16个改错）；最坏后悔排前20\%的任务占地图损失的53\%，高于翻转比例的36\%。此外，沿途测30千米在该区也显著降低了损失。纵观三个海区，一致成立的是：计算消除不了环境歧义，合理反例几乎处处存在，最坏后悔能找出损失所在，被门控标记后去观测有回报；不观测的保守改选只在黑潮显著有效，改进规则使它在其他海区无害，但没有带来显著收益。因此，门控的核心价值在于决定何时去观测，不观测时的保守改选是因海区而异的后备手段。

第六，航行器改带真实卫星海流图时，加算无效与最坏后悔风险排序两条结论依然成立，而各种响应的收益缩小且不再显著。该检验在黑潮区进行：航行器带3天前的0.25°近实时融合卫星测高地转流图，上浮取当天的同一产品，另以当天0.5°模式图作理想参照；“真实海洋”仍为再分析场；误差模型、半径与阈值在开发年份由“测高图减再分析场”标定（该产品有一次长时间中断，地图比预定日期老1天以上的快照在看到结果前即被排除）；检验计划在下载卫星数据之前写定。五项预设判定中三项成立：加算后选错率只降0.3个百分点；最坏后悔排前20\%的任务占地图损失的67\%，高于翻转比例的52\%；改进规则的点估计不高于照地图走（置信区间含0）。两项不成立：存在合理反例的任务为81.0\%，低于预设的90\%；门控后上浮取当天0.5°模式图的置信区间含0。照精算地图走的选错率为19.6\%，换成当天的测高图仍有18.5\%，换成当天0.5°模式图则降至8.6\%，说明测高图与再分析场之差主要是两种产品之间的系统差，而非时效。由此，加算无效与风险排序不依赖“粗图与真实海洋同源”的设置；主设置中门控后取保守的显著收益，则未能在真实卫星图下复现。

\begin{table}[htbp]
\centering
\caption{不同海流图、航行器速度、合理集合构造方法与海区下的适用范围与消融实验结果（测试年份；损失单位：小时；门控后上浮取图指取当天0.5°图）。}
\label{tab:scope}
\small
\begin{tabular}{@{}p{0.19\linewidth}>{\centering\arraybackslash}p{0.088\linewidth}>{\centering\arraybackslash}p{0.088\linewidth}>{\centering\arraybackslash}p{0.088\linewidth}>{\centering\arraybackslash}p{0.088\linewidth}>{\centering\arraybackslash}p{0.088\linewidth}>{\centering\arraybackslash}p{0.088\linewidth}>{\centering\arraybackslash}p{0.088\linewidth}@{}}
\toprule
设置 & 测试任务 & 存在合理反例 & 基线选错率 & 基线额外航时 & 超阈值时保守改选 & 超阈值时上浮取图 & 最大后悔值前20\%所占额外航时 \\
\midrule
黑潮主设置 & 623 & 99.7\% & 19.7\% & 0.98 & 0.70 & 0.78 & 67\% \\
粗图换成当天0.5°图 & 349 & 64.5\% & 8.0\% & 0.17 & 0.17 & 0.16 & 60\% \\
航速降至0.75米/秒 & 347 & 99.1\% & 18.4\% & 2.43 & 1.80 & 1.32 & 74\% \\
消融：合理集合改用相似日差值 & 349 & 99.4\% & 19.5\% & 0.99 & 0.95 & 0.77 & 64\% \\
湾流 & 643 & 98.8\% & 19.4\% & 3.56 & 4.00 & 1.31 & 71\% \\
阿古拉斯回流区 & 618 & 99.2\% & 27.7\% & 6.87 & 6.49 & 2.20 & 53\% \\
改带卫星测高地转流图 & 648 & 81.0\% & 19.6\% & 0.79 & 0.71 & 0.64 & 67\% \\
\bottomrule
\end{tabular}

\end{table}

\subsection{局限性讨论}
本研究存在五方面局限：第一，单次任务求解将海流快照视为定常二维场，尚未耦合约40小时航行中的流场时变；第二，主实验的旧粗图由同一再分析场平均与滞后得到，尽管第\ref{sec:scope}节改用真实的卫星测高地转流图重复了检验，但“真实海洋”仍取当天1/12°再分析场，且再分析场仍平滑了亚中尺度过程与内波扰动；第三，将绕障简化为单个圆形禁航区的左右二元决策，未扩展至多岛礁群的多分支同伦类；第四，假设选定一侧后航行器能按真实海流以最优轨迹航行，以剥离底层轨迹跟踪误差、专门评估顶层选边决策；第五，不观测的门控后取保守只在黑潮主设置下显著降低损失，在另两个海区与真实卫星图下均不显著（第\ref{sec:scope}节），此时应依靠门控触发观测来消解环境歧义。
